\documentclass[12pt,a4paper]{article}

% ============================================================
% PACKAGES
% ============================================================
\usepackage[utf8]{inputenc}
\usepackage[T1]{fontenc}
\usepackage{lmodern}
\usepackage[french]{babel}
\usepackage{geometry}
\usepackage{graphicx}
\usepackage[table]{xcolor}
\usepackage{listings}
\usepackage{booktabs}
\usepackage{array}
\usepackage{fancyhdr}
\usepackage{titlesec}
\usepackage{enumitem}
\usepackage{mdframed}
\usepackage{float}
\usepackage{microtype}
\usepackage{setspace}
\usepackage{hyperref}

% ============================================================
% MISE EN PAGE
% ============================================================
\geometry{
  headheight=25pt,
  top=2.6cm,
  bottom=2.8cm,
  left=2.8cm,
  right=2.8cm
}
\setstretch{1.12}

% ============================================================
% COULEURS
% ============================================================
\definecolor{primary}{RGB}{28, 90, 168}
\definecolor{primarylight}{RGB}{88, 145, 220}
\definecolor{secondary}{RGB}{232, 240, 252}
\definecolor{codebg}{RGB}{248, 250, 255}
\definecolor{codeframe}{RGB}{170, 195, 230}
\definecolor{success}{RGB}{34, 130, 90}
\definecolor{warning}{RGB}{200, 120, 40}
\definecolor{darkgray}{RGB}{45, 55, 72}
\definecolor{lightgray}{RGB}{95, 110, 140}
\definecolor{rowgray}{RGB}{245, 247, 250}

% ============================================================
% STYLE DU CODE
% ============================================================
\lstset{
  backgroundcolor=\color{codebg},
  frame=single,
  rulecolor=\color{codeframe},
  basicstyle=\ttfamily\footnotesize,
  keywordstyle=\color{primary}\bfseries,
  commentstyle=\color{lightgray}\itshape,
  stringstyle=\color{success},
  breaklines=true,
  showstringspaces=false,
  numbers=left,
  numberstyle=\tiny\sffamily\textcolor{lightgray},
  numbersep=6pt,
  tabsize=2,
  captionpos=b,
  aboveskip=1em,
  belowskip=1em,
  xleftmargin=0pt,
  framexleftmargin=0pt
}

% ============================================================
% EN-TETE ET PIED DE PAGE
% ============================================================
\pagestyle{fancy}
\fancyhf{}
\fancyhead[L]{\small\sffamily\textcolor{primary!85!black}{Projet M1 GIL --- Gestion de flotte}}
\fancyhead[R]{\small\sffamily\textcolor{primary!85!black}{Univ. Rouen Normandie \textperiodcentered{} 2025--2026}}
\fancyfoot[C]{\small\sffamily\textcolor{primary!70!black}{\thepage}}
\renewcommand{\headrulewidth}{0.4pt}
\renewcommand{\headrule}{\hbox to\headwidth{\color{primary!35}\leaders\hrule height \headrulewidth\hfill}}
\renewcommand{\footrulewidth}{0pt}

% ============================================================
% STYLE DES TITRES
% ============================================================
\titleformat{\section}
  {\Large\sffamily\bfseries\color{primary}}
  {\thesection.}{0.55em}{}

\titleformat{\subsection}
  {\large\sffamily\bfseries\color{primary!85!black}}
  {\thesubsection.}{0.45em}{}

\titleformat{\subsubsection}
  {\normalsize\sffamily\bfseries\color{primary!75!black}}
  {\thesubsubsection.}{0.4em}{}

\titlespacing*{\section}{0pt}{1.6ex plus .2ex}{1ex plus .1ex}
\titlespacing*{\subsection}{0pt}{1.3ex plus .15ex}{0.7ex plus .1ex}

% ============================================================
% BOITE INFO
% ============================================================
\newmdenv[
  backgroundcolor=secondary,
  linecolor=primary,
  linewidth=1.1pt,
  leftline=true,
  rightline=false,
  topline=false,
  bottomline=false,
  innerleftmargin=11pt,
  innerrightmargin=11pt,
  innertopmargin=9pt,
  innerbottommargin=9pt,
  skipabove=0.6em,
  skipbelow=0.6em
]{infobox}

\hypersetup{
  colorlinks=true,
  linkcolor=primary!90!black,
  urlcolor=primarylight,
  pdfborderstyle={/S/U/W 0}
}

% ============================================================
% DOCUMENT
% ============================================================
\begin{document}

% ------------------------------------------------------------
% PAGE DE TITRE
% ------------------------------------------------------------
\begin{titlepage}
  \centering
  \vspace*{1cm}

  {\Large\textbf{Universit\'{e} de Rouen Normandie}\\[0.3cm]
  \large Master 1 -- G\'{e}nie Informatique et Logiciel}

  \vspace{1.5cm}
  {\color{primary!55}\rule{\linewidth}{0.5mm}}\\[0.45cm]

  {\Huge\sffamily\bfseries\color{primary} Projet Gestion de Flotte\\[0.35cm]}
  {\LARGE\sffamily\bfseries\color{primary!80!black} Rapport de Semaine 6}\\[0.35cm]
  {\large\sffamily\color{primary!65!black} Architecture Frontend, S\'{e}curit\'{e} SSO \& Tests E2E}

  {\color{primary!55}\rule{\linewidth}{0.5mm}}\\[1.5cm]

  \begin{minipage}{0.45\textwidth}
    \begin{flushleft}
      \textbf{\'{E}tudiants :}\\
      Ahc\`{e}ne Amouchas \\
      Florian Pepin \\[0.5cm]
      \textbf{Encadrants :}\\
      Lydia \& Luc
    \end{flushleft}
  \end{minipage}
  \hfill
  \begin{minipage}{0.45\textwidth}
    \begin{flushright}
      \textbf{Ann\'{e}e universitaire :}\\
      2025 -- 2026 \\[0.5cm]
      \textbf{Semaine :}\\
      6 / 7
    \end{flushright}
  \end{minipage}

  \vfill
  {\large 18 avril 2026}
\end{titlepage}

\tableofcontents
\newpage

% ============================================================
\section{D\'{e}veloppement du Micro-Frontend}
% ============================================================

\subsection{Philosophie et structure en Monorepo}

L'interface utilisateur de la plateforme de gestion de flotte a \'{e}t\'{e} con\c{c}ue selon une architecture \textbf{Micro-Frontend (MFE)}, organis\'{e}e au sein d'un \textbf{Monorepo NPM Workspaces}. Cette approche permet \`{a} plusieurs d\'{e}veloppeurs de travailler en parall\`{e}le sur des domaines m\'{e}tiers distincts sans g\'{e}n\'{e}rer de conflits, tout en maintenant une coh\'{e}rence visuelle et s\'{e}curitaire globale.

L'arborescence se divise en deux cat\'{e}gories :

\begin{itemize}[leftmargin=*]
  \item \textbf{\texttt{apps/}} : applications ex\'{e}cutables de fa\c{c}on autonome (\texttt{shell}, \texttt{vehicles}, \texttt{drivers}, \texttt{maintenance}, \texttt{location}).
  \item \textbf{\texttt{packages/}} : biblioth\`{e}ques transversales partag\'{e}es (\texttt{shared-auth}, \texttt{shared-ui}, \texttt{shared-client}).
\end{itemize}

\subsection{Module Federation avec Vite.js}

L'assemblage dynamique des micro-applications s'effectue gr\^{a}ce au plugin \textbf{Module Federation} de Vite. L'application \textbf{Shell} joue le r\^{o}le d'h\^{o}te : elle d\'{e}finit les routes, applique la politique de s\'{e}curit\'{e} et charge \`{a} la demande les composants expos\'{e}s par les applications distantes (\textit{Remotes}).

\begin{lstlisting}[language=bash, caption=Configuration Module Federation --- shell/vite.config.ts]
federation({
  name: 'shell',
  remotes: {
    vehicles_app:    'http://localhost:5002/assets/remoteEntry.js',
    drivers_app:     'http://localhost:5003/assets/remoteEntry.js',
    maintenance_app: 'http://localhost:5004/assets/remoteEntry.js',
    location_app:    'http://localhost:5005/assets/remoteEntry.js',
  },
  shared: ['react', 'react-dom', '@apollo/client', 'graphql'],
})
\end{lstlisting}

\subsection{Composants React et chargement asynchrone}

Chaque module distant est import\'{e} dynamiquement via \texttt{React.lazy()} et envelopp\'{e} dans un \texttt{<Suspense>}, ce qui diff\`{e}re le t\'{e}l\'{e}chargement du bundle jusqu'au moment o\`{u} l'utilisateur acc\`{e}de \`{a} la route correspondante.

\begin{lstlisting}[language=Java, caption=shell/src/App.tsx --- chargement lazy des remotes]
const VehicleList = React.lazy(() =>
  import('vehicles_app/VehicleList').then(m => ({ default: m.default || m }))
);

<Route path="vehicles" element={
  <Suspense fallback={<div>Chargement du module vehicules...</div>}>
    <VehicleList />
  </Suspense>
} />
\end{lstlisting}

\subsection{Communication avec la Gateway GraphQL}

Un paquetage partag\'{e} \texttt{@flotte/shared-client} base sur \texttt{@apollo/client} injecte automatiquement le jeton JWT dans l'en-t\^{e}te \texttt{Authorization} de chaque requ\^{e}te via un \textit{Context Link} Apollo.

\begin{lstlisting}[language=Java, caption=packages/shared-client --- injection du token JWT]
const authLink = new ApolloLink((operation, forward) => {
  const token = keycloak.token;
  operation.setContext({
    headers: { Authorization: token ? `Bearer ${token}` : '' },
  });
  return forward(operation);
});
\end{lstlisting}

% ============================================================
\section{Int\'{e}gration Keycloak (SSO, login, logout)}
% ============================================================

\subsection{Choix de Keycloak et protocole OIDC}

L'authentification a \'{e}t\'{e} d\'{e}l\'{e}gu\'{e}e \`{a} un serveur \textbf{Keycloak} conteneuris\'{e}, impl\'{e}mentant le protocole \textbf{OAuth2 / OpenID Connect (OIDC)}. Aucun formulaire de connexion n'est pr\'{e}sent dans le code React : les requ\^{e}tes non authentifi\'{e}es sont automatiquement redirig\'{e}es vers Keycloak (\textit{Standard Flow} avec \textit{PKCE}).

\begin{infobox}
\textbf{Flux de connexion :} L'utilisateur acc\`{e}de \`{a} l'application $\rightarrow$ Keycloak intercepte et affiche sa page de login $\rightarrow$ Apr\`{e}s validation des identifiants, Keycloak retourne un \textbf{JWT} (access token + refresh token) $\rightarrow$ L'application est mont\'{e}e et les requ\^{e}tes GraphQL sont autoris\'{e}es.
\end{infobox}

\subsection{Configuration Keycloak}

Keycloak est configur\'{e} avec :
\begin{itemize}[leftmargin=*]
  \item Le \textit{realm} \texttt{gestion-flotte} contenant les utilisateurs et les r\^{o}les m\'{e}tiers.
  \item Un client OIDC \texttt{frontend} de type \textit{public} (sans secret), utilisant \texttt{onLoad: 'login-required'} pour forcer l'authentification d\`{e}s l'acc\`{e}s.
  \item Le chemin HTTP relatif \texttt{/auth} (\texttt{KC\_HTTP\_RELATIVE\_PATH=/auth}) permettant son exposition sous \texttt{flotte.local/auth} via l'ingress \textbf{Traefik} (et non Nginx), utilis\'{e} comme contr\^{o}leur d'ingress dans notre d\'{e}ploiement Helm.
\end{itemize}

\subsection{Utilisateurs pr\'{e}-configur\'{e}s}

Trois utilisateurs sont automatiquement inject\'{e}s dans le realm lors du d\'{e}marrage, couvrant les trois profils m\'{e}tiers de l'application :

\begin{table}[H]
\centering
\renewcommand{\arraystretch}{1.4}
\begin{tabular}{
  >{\ttfamily\raggedright\arraybackslash}p{3cm}
  >{\ttfamily\centering\arraybackslash}p{3cm}
  >{\raggedright\arraybackslash}p{7cm}
}
\toprule
\rowcolor{secondary}
\textbf{Identifiant} & \textbf{Mot de passe} & \textbf{R\^{o}le} \\
\midrule
admin      & admin      & Administrateur (acc\`{e}s complet) \\
\rowcolor{rowgray}
technicien & technicien & Technicien (maintenance et v\'{e}hicules) \\
manager    & manager    & Manager (tableaux de bord et conducteurs) \\
\bottomrule
\end{tabular}
\caption{Utilisateurs par d\'{e}faut du realm \texttt{gestion-flotte}}
\end{table}

\subsection{Package \texttt{@flotte/shared-auth}}

La m\'{e}canique d'int\'{e}gration Keycloak est encapsul\'{e}e dans le paquetage transversal \texttt{@flotte/shared-auth}, inject\'{e} \`{a} la racine du Shell.

\begin{enumerate}[leftmargin=*]
  \item \textbf{Verrou d'initialisation :} Un verrou m\'{e}moris\'{e} hors de la port\'{e}e du composant garantit une unique instanciation du client Keycloak, m\^{e}me en mode \textit{React Strict Mode}.
  \item \textbf{Blocage du rendu :} L'arbre DOM n'est mont\'{e} qu'apr\`{e}s r\'{e}solution compl\`{e}te de l'initialisation, \'{e}vitant tout \textit{flickering} de routes prot\'{e}g\'{e}es.
  \item \textbf{Hook \texttt{useAuth} :} Expose une fa\c{c}ade permettant \`{a} n'importe quel micro-frontend de r\'{e}cup\'{e}rer l'identit\'{e} utilisateur (\texttt{username}, \texttt{roles}) ou de d\'{e}clencher une d\'{e}connexion propre.
\end{enumerate}

\begin{lstlisting}[language=Java, caption=Utilisation du hook useAuth dans le Layout]
const { username, roles, logout } = useAuth();

<button onClick={logout}>
  <LogOut size={16} /> Quitter
</button>
\end{lstlisting}

% ============================================================
\section{Gestion JWT et RBAC (routes prot\'{e}g\'{e}es par r\^{o}le)}
% ============================================================

\subsection{Structure du JWT Keycloak}

Apr\`{e}s authentification, Keycloak \'{e}met un \textbf{JSON Web Token (JWT)} sign\'{e} contenant les informations d'identit\'{e} et les r\^{o}les de l'utilisateur. Le frontend ne stocke jamais le token en \texttt{localStorage} : il est conserv\'{e} en m\'{e}moire par le client Keycloak.js et renouvel\'{e} automatiquement avant expiration (\textit{token refresh silencieux}).

\begin{lstlisting}[language=bash, caption=Extrait d'un payload JWT Keycloak d\'{e}cod\'{e}]
{
  "sub": "uuid-utilisateur",
  "preferred_username": "florian",
  "realm_access": {
    "roles": ["fleet-manager", "offline_access"]
  },
  "exp": 1713480000
}
\end{lstlisting}

\subsection{Injection automatique dans Apollo}

Le \textit{Context Link} Apollo r\'{e}cup\`{e}re \`{a} chaque requ\^{e}te le token courant (\texttt{keycloak.token}) et l'injecte dans l'en-t\^{e}te HTTP. La Gateway GraphQL valide ensuite ce token aupr\`{e}s de Keycloak avant de laisser passer la requ\^{e}te.

\subsection{Contr\^{o}le d'acc\`{e}s bas\'{e} sur les r\^{o}les (RBAC)}

Les r\^{o}les extraits du JWT permettent de conditionner dynamiquement l'interface. Le hook \texttt{useAuth} expose un tableau \texttt{roles} consommable par n'importe quel composant pour masquer ou afficher des actions selon les habilitations.

\begin{lstlisting}[language=Java, caption=Affichage conditionnel selon le r\^{o}le]
const { roles } = useAuth();
const isManager = roles?.includes('fleet-manager');

{isManager && (
  <button onClick={() => openModal()}>
    Ajouter un vehicule
  </button>
)}
\end{lstlisting}

\begin{infobox}
\textbf{S\'{e}curit\'{e} en profondeur :} La v\'{e}rification RBAC c\^{o}t\'{e} frontend est uniquement cosmétique. La v\'{e}ritable barrière de s\'{e}curit\'{e} est appliqu\'{e}e c\^{o}t\'{e} serveur par la Gateway GraphQL, qui rejette tout token invalide ou insuffisamment habilit\'{e}.
\end{infobox}

% ============================================================
\section{Dashboards temps r\'{e}el}
% ============================================================

Chaque domaine m\'{e}tier poss\`{e}de sa propre vue d\'{e}di\'{e}e, consommant l'API GraphQL et affichant les donn\'{e}es sous forme de tableau ou de carte interactive. Les donn\'{e}es sont rafra\^{i}chies soit par requ\^{e}te (\texttt{useQuery}) soit par abonnement WebSocket (\texttt{useSubscription}) selon le module.

\subsection{Vue Tableau de bord}

La page d'accueil pr\'{e}sente une vue synth\'{e}tique de l'activit\'{e} du parc. Elle agr\`{e}ge en une seule requ\^{e}te GraphQL (\texttt{GetAllForStats}) les donn\'{e}es des trois microservices principaux et calcule c\^{o}t\'{e} client quatre indicateurs cl\'{e}s :

\begin{itemize}[leftmargin=*]
  \item \textbf{V\'{e}hicules disponibles} : v\'{e}hicules au statut \texttt{AVAILABLE}.
  \item \textbf{En maintenance} : interventions au statut \texttt{IN\_PROGRESS}.
  \item \textbf{Conducteurs actifs} : conducteurs au statut \texttt{ACTIVE} ou \texttt{ON\_TOUR}.
  \item \textbf{Maintenances pr\'{e}vues} : interventions au statut \texttt{SCHEDULED}.
\end{itemize}

Chaque indicateur est affich\'{e} dans une carte visuelle avec une ic\^{o}ne et un code couleur propre \`{a} son \'{e}tat m\'{e}tier.

\subsection{Vue V\'{e}hicules}

Le module \texttt{vehicles} affiche l'ensemble du parc automobile dans un tableau avec les colonnes : immatriculation, marque et mod\`{e}le, kilom\'{e}trage et statut. Chaque ligne propose des actions d'\'{e}dition et de suppression. Un formulaire modal permet la cr\'{e}ation d'un nouveau v\'{e}hicule (plaque, marque, mod\`{e}le, kilom\'{e}trage initial, charge utile, volume, carburant). Les statuts sont mat\'{e}rialis\'{e}s par des badges color\'{e}s (\texttt{AVAILABLE}, \texttt{ON\_DELIVERY}, \texttt{IN\_MAINTENANCE}, \texttt{OUT\_OF\_SERVICE}).

\subsection{Vue Conducteurs}

Le module \texttt{drivers} liste les conducteurs avec leur pr\'{e}nom, nom et statut courant. Comme pour les v\'{e}hicules, les enregistrements sont r\'{e}cup\'{e}r\'{e}s via une requ\^{e}te GraphQL (\texttt{GetDrivers}) et affich\'{e}s dans un tableau. Le statut (\texttt{ACTIVE}, \texttt{ON\_TOUR}, \texttt{UNAVAILABLE}) est rendu visuellement par un badge.

\subsection{Vue Maintenance}

Le module \texttt{maintenance} pr\'{e}sente l'ensemble des interventions planifi\'{e}es ou en cours sur le parc. Chaque ligne du tableau indique le v\'{e}hicule concern\'{e}, le type d'intervention, la date pr\'{e}visionnelle et le statut (\texttt{SCHEDULED}, \texttt{IN\_PROGRESS}, \texttt{COMPLETED}). Les donn\'{e}es sont charg\'{e}es via \texttt{maintenanceRecords} depuis la Gateway GraphQL.

\subsection{Vue Localisation --- Suivi GPS temps r\'{e}el}

Le module \texttt{location} affiche la position GPS des v\'{e}hicules sur une carte interactive \textbf{Leaflet}. Deux m\'{e}canismes sont combin\'{e}s :

\begin{enumerate}[leftmargin=*]
  \item \textbf{Requ\^{e}te initiale (\texttt{useQuery})} : r\'{e}cup\`{e}re la derni\`{e}re position connue au chargement de la page.
  \item \textbf{Abonnement temps r\'{e}el (\texttt{useSubscription})} : \'{e}coute un flux WebSocket GraphQL (\texttt{WATCH\_VEHICLE\_LOCATION}) et met \`{a} jour la carte \`{a} chaque nouveau signal GPS re\c{c}u.
\end{enumerate}

\begin{lstlisting}[language=Java, caption=location/src/App.tsx --- abonnement WebSocket]
const { data: realtimeData } = useSubscription(WATCH_VEHICLE_LOCATION, {
  variables: { vehicle_id: selectedVehicleId },
});

const location = realtimeData?.vehicleLocationUpdated
               || initialData?.vehicleLocation;
\end{lstlisting}

La position courante est transmise au composant \texttt{RecenterAutomatically} qui repositionne la carte sans rechargement gr\^{a}ce au hook \texttt{useMap()} de \textit{react-leaflet}. La vitesse et le cap courants sont affich\'{e}s en temps r\'{e}el dans l'en-t\^{e}te de la vue.

\subsubsection{Script de simulation}

Un script \texttt{scripts/simulate.sh} permet de simuler un v\'{e}hicule en mouvement en envoyant des coordonn\'{e}es GPS r\'{e}elles via gRPC directement vers le \texttt{location-service}. Le trajet part de Rouen et progresse en temps r\'{e}el, ce qui permet de visualiser le d\'{e}placement sur la carte Leaflet sans intervention manuelle.

\begin{lstlisting}[language=bash, caption=Lancement de la simulation GPS]
# Sur Docker Compose
./scripts/simulate.sh docker

# Sur Minikube (ouvre automatiquement un port-forward gRPC)
./scripts/simulate.sh minikube
\end{lstlisting}

% ============================================================
\section{Tests End-to-End avec Playwright}
% ============================================================

\subsection{Mise en place}

Les tests E2E sont impl\'{e}ment\'{e}s avec \textbf{Playwright} et configur\'{e}s au niveau du Monorepo (\texttt{frontend/playwright.config.ts}). Deux projets sont d\'{e}finis :

\begin{itemize}[leftmargin=*]
  \item \textbf{\texttt{setup}} : ex\'{e}cute l'authentification Keycloak et sauvegarde la session dans \texttt{e2e/.auth/user.json}.
  \item \textbf{\texttt{chromium}} : charge la session sauvegard\'{e}e (\texttt{storageState}) et ex\'{e}cute les tests m\'{e}tiers. Il d\'{e}pend du projet \texttt{setup}.
\end{itemize}

\begin{lstlisting}[language=Java, caption=playwright.config.ts --- configuration des projets]
projects: [
  { name: 'setup', testMatch: '**/auth.setup.ts' },
  {
    name: 'chromium',
    use: {
      ...devices['Desktop Chrome'],
      storageState: 'e2e/.auth/user.json',
    },
    dependencies: ['setup'],
  },
]
\end{lstlisting}

\subsection{Authentification automatis\'{e}e}

Le fichier \texttt{auth.setup.ts} simule la connexion Keycloak une seule fois avant la suite de tests. La session est r\'{e}utilis\'{e}e pour tous les tests suivants, \'{e}vitant une reconnexion \`{a} chaque sc\'{e}nario.

\begin{lstlisting}[language=Java, caption=e2e/tests/auth.setup.ts]
setup('authenticate via Keycloak', async ({ page }) => {
  await page.goto('/');
  await page.waitForURL(/\/auth\/realms\/gestion-flotte/);
  await page.fill('#username', 'test-user');
  await page.fill('#password', 'password');
  await page.click('#kc-login');
  await page.waitForURL('/');
  await expect(
    page.getByRole('heading', { name: 'Tableau de bord' })
  ).toBeVisible();
  await page.context().storageState({ path: authFile });
});
\end{lstlisting}

\subsection{Lancement des tests}

Un script shell \texttt{scripts/test-e2e.sh} permet de lancer les tests selon l'environnement cible. Il v\'{e}rifie la disponibilit\'{e} de l'application avant de d\'{e}marrer Playwright et ouvre automatiquement le rapport HTML en fin d'ex\'{e}cution.

\begin{lstlisting}[language=bash, caption=Lancement des tests E2E]
# Mode Docker Compose (http://localhost:3005)
./scripts/test-e2e.sh docker

# Mode Minikube (http://flotte.local)
./scripts/test-e2e.sh minikube
\end{lstlisting}

\subsection{R\'{e}capitulatif des tests}

\begin{table}[H]
\centering
\renewcommand{\arraystretch}{1.45}
\begin{tabular}{
  >{\raggedright\arraybackslash}p{3cm}
  >{\raggedright\arraybackslash}p{6.2cm}
  >{\centering\arraybackslash}p{2.2cm}
}
\toprule
\rowcolor{secondary}
\textbf{Module} & \textbf{Sc\'{e}nario test\'{e}} & \textbf{R\'{e}sultat} \\
\midrule
Authentification & Redirection vers Keycloak et connexion & \textcolor{success}{\textbf{Pass}} \\
\rowcolor{rowgray}
Authentification & Sauvegarde de session (\texttt{storageState}) & \textcolor{success}{\textbf{Pass}} \\
Dashboard & Affichage du titre \og Tableau de bord \fg{} & \textcolor{success}{\textbf{Pass}} \\
\rowcolor{rowgray}
Dashboard & Pr\'{e}sence des 4 cartes de statistiques & \textcolor{success}{\textbf{Pass}} \\
Dashboard & Chargement des valeurs sans erreur GraphQL & \textcolor{success}{\textbf{Pass}} \\
\rowcolor{rowgray}
V\'{e}hicules & Affichage de la liste des v\'{e}hicules & \textcolor{success}{\textbf{Pass}} \\
V\'{e}hicules & Cr\'{e}ation d'un v\'{e}hicule via le formulaire & \textcolor{success}{\textbf{Pass}} \\
\rowcolor{rowgray}
V\'{e}hicules & Suppression d'un v\'{e}hicule avec confirmation & \textcolor{success}{\textbf{Pass}} \\
Conducteurs & Affichage de la liste des conducteurs & \textcolor{success}{\textbf{Pass}} \\
\rowcolor{rowgray}
Conducteurs & Absence d'erreur GraphQL & \textcolor{success}{\textbf{Pass}} \\
Maintenance & Affichage de la liste des interventions & \textcolor{success}{\textbf{Pass}} \\
\rowcolor{rowgray}
Maintenance & Absence d'erreur GraphQL & \textcolor{success}{\textbf{Pass}} \\
Localisation & Chargement de la carte Leaflet & \textcolor{success}{\textbf{Pass}} \\
\rowcolor{rowgray}
Localisation & Absence d'erreur d'initialisation & \textcolor{success}{\textbf{Pass}} \\
\bottomrule
\end{tabular}
\caption{R\'{e}capitulatif des tests end-to-end Playwright}
\end{table}

% ============================================================
\newpage
\section{Bilan et perspectives pour la Semaine 7}
% ============================================================

\subsection{Travaux r\'{e}alis\'{e}s}

L'ensemble du \textbf{frontend} a \'{e}t\'{e} d\'{e}velopp\'{e} et est fonctionnel : architecture Micro-Frontend, int\'{e}gration Keycloak (SSO, JWT, RBAC), dashboards temps r\'{e}el (v\'{e}hicules, conducteurs, maintenance, localisation GPS) et tests E2E Playwright. Le p\'{e}rim\`{e}tre frontal pr\'{e}vu pour la semaine 6 est donc compl\`{e}tement trait\'{e}.

\subsection{Ce qui est pr\'{e}vu en semaine 7}

La semaine 7 comportera plusieurs travaux. Parmi ceux-ci, le \textbf{syst\`{e}me d'alertes} n'a pas encore \'{e}t\'{e} impl\'{e}ment\'{e} et figurera au programme : ce composant permettra de notifier les op\'{e}rateurs en cas d'\'{e}v\'{e}nements m\'{e}tiers critiques (v\'{e}hicule en panne, maintenance urgente, etc.).

\end{document}
